\section{Performance Characteristics}
\label{sec:performance}

All numbers in this section are measured on the new final Lux
network's reference validator, using \texttt{go test -bench}. Sources
of raw data:

\begin{itemize}[leftmargin=2em,nosep]
\item \texttt{/tmp/zap-vs-codec/} -- ZAP vs legacy \texttt{linearcodec}
  per-operation microbenchmarks; raw outputs in
  \texttt{raw\_bench\_\{m1max,evo,dbc,spark\}.txt}.
\item \texttt{/tmp/multi-validator/} -- distributed end-to-end cert
  benches at $10$, $100$, $500$ validators across three host hardware
  classes; \texttt{result-\{10,100,500\}v-\{aurora,polaris,hybrid\}.json}.
\item \texttt{/tmp/spark-gpu-bench/} -- DGX Spark GB10 GPU baseline
  for the Corona/Pulsar lattice paths; \texttt{MEDIANS.txt} and
  \texttt{raw\_bench\_spark.txt}.
\item \texttt{lux/threshold/protocols/parity/} -- Corona/Pulsar parity
  audit (\texttt{parity\_test.go}); the result that locks
  Pulsar-first in LP-217's PQ-fast tier.
\end{itemize}

\subsection{Per-operation wins (measured)}

\begin{table}[h]
\centering
\small
\begin{tabular}{@{}lrrr@{}}
\toprule
\textbf{Operation} & \textbf{Legacy (linearcodec)} & \textbf{ZAP Stack} & \textbf{Win} \\
\midrule
Parse small tx & \SI{129}{\nano\second}/op, 2 allocs & \SI{36}{\nano\second}/op, 1 alloc & $3.6\times$ \\
Parse complex tx & \SI{1008}{\nano\second}/op, 6 allocs & \SI{75}{\nano\second}/op, 1 alloc & $13.4\times$ \\
Mempool insert & parse $\to$ struct $\to$ remarshal & parse $\to$ buffer ref & $\sim 5\times$ \\
State write & unmarshal $\to$ mutate $\to$ marshal & wrap $\to$ mutate $\to$ write buffer & $\sim 3\times$ \\
Snapshot stream & re-encode for backup & direct buffer stream & $\sim 10\times$ \\
Corona aggregate $n{=}64$ (incorrect serial baseline) & \SI{14}{\second} & \SI{356}{\milli\second}--\SI{440}{\milli\second} parallel & $-97\%$ \\
\bottomrule
\end{tabular}
\caption{Per-operation cost: ZAP Stack vs legacy. The Corona row
deserves a note (Section~\ref{ssec:corona-row}).}
\label{tab:perops}
\end{table}

\subsection{End-to-end distributed-cluster cert finality}

The multi-validator bench (\texttt{/tmp/multi-validator/}) measures
end-to-end QuasarCert finality across a three-host cluster at
$10$/$100$/$500$ validator scales. The hosts span the three production
hardware families: Apple Silicon (\texttt{ra}, M1 Max), NVIDIA
Blackwell-class GPU server (\texttt{spark}, GB10), and x86\_64 Linux
(\texttt{evo}, Ryzen AI MAX+).

\begin{table}[h]
\centering
\small
\begin{tabular}{@{}llrrrr@{}}
\toprule
\textbf{Scale} & \textbf{Profile} & \textbf{Blocks/s} & \textbf{BLS leg} & \textbf{Pulsar leg} & \textbf{Corona leg} \\
\midrule
$10$v   & Hybrid-1s & 0.88 & \SI{159}{\milli\second} & \SI{203}{\milli\second} & \SI{1135}{\milli\second} \\
$100$v  & Hybrid-1s & 0.39 & \SI{162}{\milli\second} & \SI{298}{\milli\second} & \SI{2593}{\milli\second} \\
$500$v  & Hybrid-1s & 0.38 & \SI{187}{\milli\second} & \SI{408}{\milli\second} & \SI{2655}{\milli\second} \\
\bottomrule
\end{tabular}
\caption{Leg arrival $p_{50}$ in milliseconds on the Hybrid-1s
profile. CPU baseline; Corona GPU acceleration pending. Source:
\texttt{/tmp/multi-validator/result-*v-hybrid.json}.}
\label{tab:e2e}
\end{table}

At $10$v Hybrid-1s the BLS preliminary-finality lands at
\SI{159}{\milli\second}; Pulsar lands at \SI{203}{\milli\second};
Corona lands at \SI{1.135}{\second}. At $100$v and $500$v Corona's CPU
cost dominates; the architectural shape (LP-217 PQ-fast = BLS +
Pulsar, PQ-strict adds Corona on top) is precisely the mode-tier
degradation pattern that lets relying parties at PQ-fast act in
\SI{200}{\milli\second} while PQ-strict relying parties wait the full
$\sim$\SI{2.6}{\second}.

\subsection{The Corona row deserves a note}
\label{ssec:corona-row}

The legacy serial-aggregate Corona implementation that yielded
\SI{14}{\second} per round at $n{=}64$ was incorrect, not a baseline.
The parallel implementation, which is what ships, lands at
\SI{356}{\milli\second}--\SI{440}{\milli\second}. The win in the table
is real and reflects the correct implementation cost; it is not a
straw-man comparison.

\subsection{Parity audit: Pulsar is $141\times$ cheaper than Corona at $n{=}64$}

The Corona/Pulsar parity audit
(\texttt{lux/threshold/protocols/parity/parity\_test.go}) measured
the two schemes apples-to-apples on Apple M1 Max + DGX Spark GB10 and
found Corona is \emph{$141\times$ slower than Pulsar at $n{=}64$
validators} (\SI{3.2}{\second}--\SI{4.0}{\second} vs
\SI{23}{\milli\second}--\SI{28}{\milli\second} total signing
round-trip). This is not an optimization gap -- it is algorithmic.
Pulsar is a reconstruct-on-aggregator Shamir scheme over the
unmodified FIPS 204 ML-DSA sign; Corona is a genuine 2-round
threshold lattice scheme (Corona family,). The
result locks the LP-217 PQ-fast tier at \texttt{BLS + Pulsar}; Corona
is added on top at PQ-strict for the never-trust-aggregator security
property.
